
\section{Вывод уравнений обратной задачи кинематики для платформ с кривошипно-шатунными опорами} \label{sec:InverseKinematics}

\subsection{ОЗК для трехстепенной СП}

Кинематический анализ для платформ с тремя степенями свободы широко известен как основное средство для решения задачи о базовом управлении платформой путём преобразования координат \cite{Li2005,Ng2006,Lukanin2005}.
\label{subsec:InvKin3DOFMotionPlatform}
Кинематическая схема СП 3DOF показана на рисунке \ref{img:3dof_kinematics_scheme}.
~Введены следующие обозначения:
$R$ -- длина кривошипа, $L$ -- длина шатуна, $a$ -- расстояние по оси X до двигателей 1 и 2, $b$ - расстояние по оси Y до двигателей 1 и 2, $c$ -- расстояние по оси X до двигателя 3, $\alpha$ -- угол крена, $\beta$ -- угол тангажа, $z_p = \left(O O_1\right)$ -- подъем центра подвижной платформы, $h = z_p - L$ -- смещение центра по высоте относительно среднего значения $z_p$, $q_1, q_2, q_3$ -- углы поворота кривошипов.

\begin{figure}[h]
	\center
	\includegraphics[width=0.75\textwidth]{3dof_kinematics_scheme}
	\caption{Кинематическая схема подвижности 3DOF}
	\label{img:3dof_kinematics_scheme}
\end{figure}

В ходе решения ОЗК необходимо определить вектор $\mathbf{q}=(q_1^*, q_2^*, q_3^*)^T$ обобщенных коодинат~по~заданному вектору~положения~платформы $\mathbf{X}=(\alpha,~\beta,~z_p^*)^T$. Для вывода уравнений ОЗК выразим координаты точек С1, С2, С3 через углы крена, тангажа и подъем. Для этого запишем матрицы поворота относительно осей $x$ и $y$

\begin{equation}
	\label{eq:KinematicsRotationMatrices}
	\begin{aligned}
	& R_x =
	\left(
		\begin{array}{ccc}
			1 &	0 				 &	0 \\
			0 &	\cos\alpha   &	-\sin\alpha\\
			0 &	\sin\alpha	 &	\cos\alpha
		\end{array}
	\right),
	\\
	& Ry =
	\left(
		\begin{array}{ccc}
			\cos\beta &	0   &	\sin\beta \\
			0              &	1   & 	0\\
			-\sin\beta &	0	& 	\cos\beta
		\end{array}
	\right).
	\end{aligned}
\end{equation}

Тогда координаты точек крепления шатунов к подвижной платформе:
\begin{equation}
\label{eq:KinematicsUpperPoints}
	\begin{aligned}
		\left(\begin{array}{ccc} x_{C1} \\ y_{C1} \\ z_{C1} \end{array}\right) = R_x R_y
		\left(\begin{array}{ccc} -a \\ -b \\ 0	\end{array}\right) +
		\left(\begin{array}{ccc} 0 \\ 0 \\ z_p^* \end{array}\right) =
		\left(\begin{array}{lr}
				a\cos{\beta} \\
				-b\cos{\alpha}-a\sin{\alpha}\sin{\beta} \\
				z_p^* - b\sin{\alpha} + a\cos{\alpha}\sin{\beta}\end{array}\right)\\
		\left(\begin{array}{ccc} x_{C2} \\ y_{C2} \\ z_{C2} \end{array}\right) = R_x R_y
		\left(\begin{array}{ccc} -a \\ b \\ 0	\end{array}\right) +
		\left(\begin{array}{ccc} 0 \\ 0 \\ z_p^* \end{array}\right) =
		\left(\begin{array}{lr}
				a\cos{\beta} \\
				b\cos{\alpha}-a\sin{\alpha}\sin{\beta} \\
				z_p^* + b\sin{\alpha} + a\cos{\alpha}\sin{\beta}\end{array}\right)\\
		\left(\begin{array}{ccc} x_{C3} \\ y_{C3} \\ z_{C3} \end{array}\right) = R_x R_y
		\left(\begin{array}{ccc} c \\ 0 \\ 0	\end{array}\right) +
		\left(\begin{array}{ccc} 0 \\ 0 \\ z_p^* \end{array}\right) =
		\left(\begin{array}{lr}
				c\cos{\beta} \\
				c\sin{\alpha}\sin{\beta} \\
				z_p^* - c\cos{\alpha}\sin{\beta}\end{array}\right)
	\end{aligned}
\end{equation}

Т.к. кривошипы вращаются строго в плоскости $yOz$, справедливы следующие соотношения:
\begin{equation}
\label{eq:KinematicsCrankPoints}
	\begin{aligned}
		& \left(\begin{array}{lr} x_{B1} \\ y_{B1} \\ z_{B1} \end{array}\right) =
		\left(\begin{array}{lr} -a \\ -b + R\sin{q_1^*} \\ -R\cos{q_1^*} \end{array}\right);\\
		& \left(\begin{array}{lr} x_{B2} \\ y_{B2} \\ z_{B2} \end{array}\right) =
		\left(\begin{array}{lr} -a \\ b + R\sin{q_2^*} \\ -R\cos{q_2^*} \end{array}\right);\\
		& \left(\begin{array}{lr} x_{B3} \\ y_{B3} \\ z_{B3} \end{array}\right) =
		\left(\begin{array}{lr} c \\ R\sin{q_3^*} \\ -R\cos{q_3^*} \end{array}\right).
	\end{aligned}
\end{equation}

Далее, т.к. расстояние между точками $B_i$ и $C_i$, $i=1\ldots3$ неизменно и равно $L$, можно записать следующие уравнения:
\begin{equation}
\label{eq:KinematicsRodLengths}
	\begin{aligned}
		 (x_{B1} - x_{C1})^2 + (y_{B1} - y_{C1})^2 + (z_{B1} - z_{C1})^2 = L^2;\\
		 (x_{B2} - x_{C2})^2 + (y_{B2} - y_{C2})^2 + (z_{B2} - z_{C2})^2 = L^2;\\
		 (x_{B3} - x_{C3})^2 + (y_{B3} - y_{C3})^2 + (z_{B3} - z_{C3})^2 = L^2.
	\end{aligned}
\end{equation}

После подстановки \ref{eq:KinematicsCrankPoints} и \ref{eq:KinematicsUpperPoints} в \ref{eq:KinematicsRodLengths} и дополнительных преобразований получаем следующие уравнения:

\begin{equation}
\label{eq:KinematicsDriveAnglesFromPosEq}
\begin{aligned}
\bar{A_i} + \bar{B_i}\sin{q_i} + \bar{C_i}\cos{q_i} = 0,~~i = 1,2,3
\end{aligned}
\end{equation}

где коэффициенты $\bar{A_i}, \bar{B_i}, \bar{C_i}$ зависят от $\alpha, \beta, z_p$ в соответствии с формулами \ref{eq:KinematicsUpperPoints},\ref{eq:KinematicsCrankPoints}.

\begin{equation}
\label{eq:KinematicsSquareEqKoefficients}
\begin{aligned}
	& \bar{A_i} = (x_{Ci} - x_{Ai})^2 + (y_{Ci} - y_{Ai})^2 + (z_{Ci} - z_{Ai})^2 + R^2 - L^2,\\
	& \bar{B_i} = -2R(y_{Ci} - y_{Ai}),\\
	& \bar{C_i} = 2R(z_{Ci} - z_{Ai}).
\end{aligned}
\end{equation}

Перейдя от $\sin{q_i}$ и $\cos{q_i}$ к половинному углу через $\tg\left({q_i/2}\right)$, получим следующие квадратные уравнения:
\begin{equation}
\label{eq:KinematicsSquareEquations}
\begin{aligned}
	\left(\bar{A_i}-\bar{C_i}\right)\tg^2{\frac{q_i}{2}} + 2B_i\tg{\frac{q_i}{2}} + \left(\bar{A_i}+\bar{C_i}\right)=0
\end{aligned}
\end{equation}
Из решения уравнения \ref{eq:KinematicsSquareEquations} получаем два корня:

\begin{equation}
\label{eq:KinematicsPhiRoots}
\begin{aligned}
	q_i = -2\arctan\frac{\bar{B_i}\pm\sqrt{\bar{B_i^2} + \bar{C_i^2} - \bar{A_i^2}}}{\bar{A_i} - \bar{C_i}}
\end{aligned}
\end{equation}

Оба корня с точки зрения кинематики платформы идентичны. Выбором знака <<+>> или <<-->> в уравнении \ref{eq:KinematicsPhiRoots} можно задать область значений обобщенной координаты $(q_i < 0)$ или $(q_i > 0)$, т.е. определить соответствующую рабочую полуплоскость вращения кривошипа. Как будет показано далее, выбор рабочей плоскости сказывается на динамических характеристиках СП.

\subsection{Аппроксимация уравнений ПЗК}

Для дальнейшего решения и анализа уравнений динамики СП численного решения ПЗК оказывается недостаточно, вследствие чего требуется также получить аналитические зависимости $\mathbf{X}=f\left(\mathbf{q}\right)$.
Как известно, решение ПЗК представляет сложность из-за наличия избыточных кинематических связей \cite{2014c,Bonev2000,Niu2013,Lapikov2014,2007,Charters2009,Kucuk2006,Hu2011,Jiang2008}.

Путем подстановки \ref{eq:KinematicsCrankPoints} и \ref{eq:KinematicsUpperPoints} в \ref{eq:KinematicsRodLengths} получим систему уравнений ПЗК:
\begin{equation}
\label{eq:DirectKinematicsProblemFormulation}
\begin{aligned}
& F_1\left(X\right)=\left(b\sin{\alpha}\sin{\beta}-a\cos{\beta}+a\right)^2+\left(b\cos{\alpha}+b-R\sin{q_1}\right)^2+\\
&+\left(z_p+a\sin{\beta}+b\cos{\beta}\sin{\alpha}+R\cos{q_1}\right)^2-L^2=0;\\
& F_2\left(X\right)=\left(-b\sin{\alpha}\sin{\beta}-a\cos{\beta}+a\right)^2+\left(b\cos{\alpha}+b+R\sin{q_2}\right)^2+\\
&+\left(z_p+a\sin{\beta}-b\cos{\beta}\sin{\alpha}+R\cos{q_2}\right)^2-L^2=0;\\
& F_3\left(X\right)=\left(c\cos{\beta}-c\right)^2+\left(R\sin{q_3}\right)^2+\left(z_p-c\sin{\beta}+R\cos{q_3}\right)^2-L^2=0.
\end{aligned}
\end{equation}

Т.к. уравнения системы \ref{eq:DirectKinematicsProblemFormulation} напрямую не обращаются, необходимо упростить выражения, в которые входят тригонометрические и степенные функции от $\alpha,~\beta,~h$. Для этого разложим $\alpha, \beta, h$ уравнения  в ряд по степеням $\kappa = R/L$:
\begin{equation}
\label{eq:DirectKinematicsSeriesExpansion}
\begin{aligned}
	\alpha = C_{1,\alpha}\kappa + C_{2,\alpha}\kappa^2, \\
	\beta = C_{1,\beta}\kappa + C_{2,\beta}\kappa^2, \\
	h = C_{1,h}\kappa + C_{2,h}\kappa^2.
\end{aligned}
\end{equation}

Для разложения выбрана величина $R/L$, т.к. вследствие особенностей проектирования КШМ данное соотношение всегда меньше единицы. Для исследуемого стенда $R = 0{,}15$~м, $L = 0{,}56$~м, то есть $\kappa = 0{,}268$, $\kappa^2 = 0{,}072$, $\kappa^3 = 0{,}019$. Следовательно, высокими порядками $\kappa$ можно пренебречь без существенных потерь точности вычислений $\alpha,~\beta,~h$.

Ниже выписаны коэффициенты первого и второго порядка; расчётные реализации (как MATLAB-, так и Python-версия) удерживают и третий порядок, коэффициенты которого приведены в \ref{eq:DirectKinematicsCoefsApprox3}.

Для получения коэффициентов $C_{1,\alpha},C_{2,\alpha},C_{1,\beta},C_{2,\beta}$ раскроем уравнения длин шатунов \ref{eq:KinematicsRodLengths}:
\begin{equation}
\label{eq:DirectKinematicsRodLengthsFull}
	\begin{aligned}
	0 = & h^2 + 2hL + 2a^2\left(1-\cos\beta\right) + 2b^2\left(1-\cos\alpha\right) + R^2 + 2Rb\sin{q_1}\left(1 - \right.\\
	&\left.-\cos\alpha\right) + 2\left(h+L\right)R\cos{q_1} + 2\left(h+L\right)\left(a\sin\beta + b\sin\alpha\cos\beta\right)+\\
	& + 2b\sin\alpha\left(a\sin\beta + R\cos{q_1}\cos\beta\right) + 2aR\cos{q_1}\sin\beta \\
	0 = & h^2 + 2hL + 2a^2\left(1-\cos\beta\right) + 2b^2\left(1-\cos\alpha\right) + R^2 - 2Rb\sin{q_2}\left(1 -\right.\\
	&\left. - \cos\alpha\right) + 2\left(h+L\right)R\cos{q_2} + 2\left(h+L\right)\left(a\sin\beta - b\sin\alpha\cos\beta\right)-\\
	& - 2b\sin\alpha\left(a\sin\beta + R\cos{q_2}\cos\beta\right) + 2aR\cos{q_2}\sin\beta\\
	0 = & h^2 + 2hL + 2c^2\left(1-\cos\beta\right) + R^2 + 2\left(h+L\right)R\cos{q_3}-\\
	& - 2\left(h+L\right)c\sin\beta - 2cR\cos{q_3}\sin\beta.
	\end{aligned}
\end{equation}

Для тригонометрических функций применим разложение Тейлора того же порядка:
\begin{equation}
\label{eq:DirectKinematicsSinCosSeries}
\begin{aligned}
	& \sin\alpha \approx \alpha-\frac{\alpha^3}{6} \approx C_{1,\alpha}\kappa + C_{2,\alpha}\kappa^2,\\
	& \cos\alpha \approx 1 - \frac{\alpha^2}{2} \approx 1 - \frac{C^2_{1,\alpha}}{2}\kappa^2,\\
	& \sin\beta \approx \beta-\frac{\beta^3}{6} \approx C_{1,\beta}\kappa + C_{2,\beta}\kappa^2,\\
	& \cos\beta \approx 1 - \frac{\beta^2}{2} \approx 1 - \frac{C^2_{1,\beta}}{2}\kappa^2.
\end{aligned}
\end{equation}

В результате решения полиномиальной относительно $\kappa$ системы получим следующие коэффициенты разложения функций $\alpha,\beta,h$.

Коэффициенты первого порядка при $\kappa$:
\begin{equation}
\label{eq:DirectKinematicsCoefsApprox1}
\begin{aligned}
	& C_{1,\alpha} = \frac{L}{2b}\left(\cos{q_2} - \cos{q_1}\right),\\
	& C_{1,\beta} = \frac{L}{a+c}\left(\cos{q_3} - \frac{1}{2}\left(\cos{q_1} + \cos{q_2}\right)\right),\\
	& C_{1,h} = -\frac{L}{2\left(a+c\right)}\left(2a\cos{q_3} + c\left(\cos{q_1} + \cos{q_2}\right)\right).
\end{aligned}
\end{equation}

Коэффициенты второго порядка при $\kappa^2$:
\begin{equation}
\label{eq:DirectKinematicsCoefsApprox2}
\begin{aligned}
	C_{2,\alpha} & = - \frac{1}{2}\left(\cos{q_1} + \cos{q_2}\right)C_{1,\alpha}=\frac{L}{4b}\left(\cos^2{q_1} - \cos^2{q_2}\right)\\
	C_{2,\beta} & = -\frac{L}{2\left(a+c\right)}\left[\cos^2{q_3} - \frac{1}{2}\left(\cos^2{q_1} + \cos^2{q_2}\right)\right],\\
	C_{2,h} & = cC_{2,\beta} + \frac{L}{2}\left(\cos^2{q_3} - 1\right) = \\
	& = \frac{L}{2\left(a+c\right)}\left[a\cos^2{q_3} + \frac{c}{2}\left(\cos^2{q_1} + \cos^2{q_2}\right)\right] - \frac{L}{2}.
\end{aligned}
\end{equation}

Коэффициенты третьего порядка при $\kappa^3$, удерживаемые в расчётных реализациях:
\begin{equation}
\label{eq:DirectKinematicsCoefsApprox3}
\begin{aligned}
	C_{3,\alpha} & = -\frac{1}{4}\left(\sin{q_1} + \sin{q_2}\right)C^2_{1,\alpha} + \frac{1}{6}C^3_{1,\alpha} + \frac{1}{2}C_{1,\alpha}C^2_{1,\beta},\\
	C_{3,\beta} & = \frac{b}{4\left(a+c\right)}\left(\sin{q_2} - \sin{q_1}\right)C^2_{1,\alpha} + \frac{1}{6}C^3_{1,\beta},\\
	C_{3,h} & = \frac{bc}{4\left(a+c\right)}\left(\sin{q_2} - \sin{q_1}\right)C^2_{1,\alpha}.
\end{aligned}
\end{equation}

Окончательно, подставив \ref{eq:DirectKinematicsCoefsApprox1} и \ref{eq:DirectKinematicsCoefsApprox2} в \ref{eq:DirectKinematicsSeriesExpansion} получим следующие формулы для $\alpha,~\beta,~h$:
\begin{equation}
\label{eq:DireKinematicsAnalytic}
\begin{aligned}
	& h \approx \frac{-R}{1+c/a}\left[\cos{q_3} + \frac{c}{2a}\left(\cos{q_1} + \cos{q_2}\right) + \frac{R}{2L}\left(\sin^2{q_3} + \right.\right.\\
	& \left.\left. + \frac{c}{2a}\left(\sin^2{q_1} + \sin^2{q_2}\right)\right)\right], \\
	& \alpha \approx \frac{R}{2b}\left[\cos{q_2}-\cos{q_1} + \frac{R}{2L}\left(\cos^2{q_1}-\cos^2{q_2}\right)\right],\\
	& \beta \approx \frac{R}{a+c}\left[\cos{q_3} - \frac{1}{2}\left(\cos{q_1} + \cos{q_2}\right)\right].
\end{aligned}
\end{equation}

\section{Расчет векторов шатунов}
Для решения задач динамики необходимо определять положение точек шатунов в каждый момент времени.

Пусть заданы углы $q_1,~q_2,~q_3$. Векторы вдоль шатунов от кривошипов к платформе обозначим $S_1,~S_2,~S_3$.

С точностью до третьего порядка относительно $\kappa$ их компоненты определяются следующим образом.

Первый шатун:

\begin{equation}
\begin{aligned}
\frac{1}{L}S_1 &= \begin{pmatrix}0 \\ 0 \\ 1\end{pmatrix} + \kappa\begin{pmatrix}0 \\ -\sin{q_1}\\ 0\end{pmatrix} + \kappa^2\begin{pmatrix}\frac{a}{2L}C^2_{1,\beta} + \frac{b}{L}C_{1,\alpha}^2\\-\frac{b}{2L}C_{1,\alpha}^2\\ -\frac{1}{2}\sin^2{q_1} \\ \end{pmatrix} + \\
& + \kappa^3\frac{1}{L}\begin{pmatrix}aC_{1,\beta}C_{2,\beta} + bC_{1,\alpha}C_{2,\beta} + bC_{2,\alpha}C_{1,\beta} \\ -bC_{1,\alpha}C_{2,\alpha} \\ -\frac{b}{2}\sin{q_1}C_{1,\alpha}^2\end{pmatrix}
\end{aligned}
\end{equation}

Для второго шатуна:

\begin{equation}
\begin{aligned}
\frac{1}{L}S_2 &= \begin{pmatrix}0 \\ 0 \\ 1\end{pmatrix} + \kappa\begin{pmatrix}0 \\ -\sin{q_2}\\ 0\end{pmatrix} + \kappa^2\begin{pmatrix}\frac{a}{2L}C^2_{1,\beta} + \frac{b}{L}C_{1,\alpha}^2\\\frac{b}{L}C_{1,\alpha}C_{1,\beta} -\frac{b}{2L}C_{1,\alpha}^2\\ -\frac{1}{2}\sin^2{q_2} \\ \end{pmatrix} + \\
& + \kappa^3\frac{1}{L}\begin{pmatrix}aC_{1,\beta}C_{2,\beta} + bC_{1,\alpha}C_{2,\beta} + bC_{2,\alpha}C_{1,\beta} \\ -bC_{1,\alpha}C_{2,\alpha} \\ +\frac{b}{2}\sin{q_2}C_{1,\alpha}^2\end{pmatrix}
\end{aligned}
\end{equation}

Третий шатун:

\begin{equation}
\begin{aligned}
	\frac{1}{L}S_3 = \begin{pmatrix}0 \\ 0 \\ 1\end{pmatrix} + \kappa\begin{pmatrix}0 \\ -\sin{q_3}\\ 0\end{pmatrix} + \kappa^2\begin{pmatrix}-\frac{c}{2L}C^2_{1,\beta}\\0\\-\frac{1}{2}\sin^2{q_3}\end{pmatrix} + \kappa^3\frac{1}{L}\begin{pmatrix}-\frac{c}{L}C_{1,\beta}C_{2,\beta}\\0\\0\end{pmatrix}
\end{aligned}
\end{equation}

\subsection{Расчет векторов шатунов в системе координат платформы}

Для дальнейших расчетов также потребуется вычисление координат точек шатунов в подвижной системе координат.

Пересчет координат осуществляется по формуле:
\begin{equation}
\begin{aligned}
	r^0 = R_xR_y\left(r - e_3\left(L + h\right)\right)
\end{aligned}
\end{equation}

Для сокращения последующих выкладок введем обозначения для модифицированных синусов:

\begin{equation}
\begin{aligned}
	sq_i = \sin{q_i} - C_{1,\alpha},\quad i =1,\ldots,3
\end{aligned}
\end{equation}

для показателей крена и тангажа:

\begin{equation}
\begin{aligned}
	\mu & = \frac{1}{4}C_{1,\alpha}\left(\sin{q_2} - \sin{q_1}\right) = \frac{L}{8b}\left(\cos{q_2}-\cos{q_1}\right)\left(\sin{q_2}-\sin{q_1}\right),\\
	\lambda & = \frac{1}{4}\left(2\sin{q_3}-\sin{q_1}-\sin{q_2}\right)C_{1,\alpha}
\end{aligned}
\end{equation},

скорректированного второго коэффициента $C_{2,\beta}$:

\begin{equation}
\begin{aligned}
	\gamma_1 & = \frac{a}{2L}C^2_{1,\beta} - C_{2,\beta} + \frac{b}{L}C_{1,\alpha}C_{1,\beta}, \\
	\gamma_2 & = \frac{a}{2L}C^2_{1,\beta} - C_{2,\beta} - \frac{b}{L}C_{1,\alpha}C_{1,\beta}, \\
	\gamma_3 & = -\frac{c}{2L}C^2_{1,\beta} - C_{2,\beta}.
\end{aligned}
\end{equation}

для скорректированного второго коэффициента -- $C_{2,\alpha}$:

\begin{equation}
\begin{aligned}
	\delta_1 = C_{2,\alpha} - \frac{b}{2L}C^2_{1,\alpha},\quad \delta_2 = C_{2,\alpha} + \frac{b}{2L}C^2_{1,\alpha},\quad \delta_3 = C_{2,\alpha}.
\end{aligned}
\end{equation}

Тогда для первого шатуна:

\begin{equation}
\begin{aligned}
	\frac{1}{L}S_1^0 = \begin{pmatrix} 0 \\ 0 \\ 1\end{pmatrix} + \kappa\begin{pmatrix}-C_{1,\beta} \\ -sq_1 \\ 0 \end{pmatrix} + \kappa^2\begin{pmatrix}\gamma_1 \\ \delta_1 \\ -\frac{1}{2}\left(C^2_{1,\beta} + sq^2_1\right)\end{pmatrix} + \\ \kappa^3\begin{pmatrix}
	-\mu \frac{b}{a+c}C_{1,\alpha} + \frac{b}{L}C_{1,\alpha}C_{2,\beta} - \frac{1}{L}C_{1,\beta}C_{2,h}\\
	\frac{a}{L}C_{1,\alpha}C_{2,\beta} + \frac{1}{L}C_{1,\alpha}C_{2,h} - \mu C_{1,\alpha} \\
	C_{1,\beta}\gamma_1 + sq_1\delta_1
	\end{pmatrix}
\end{aligned}
\end{equation}

Для второго шатуна:

\begin{equation}
\begin{aligned}
	\frac{1}{L}S_2^0 = \begin{pmatrix} 0 \\ 0 \\ 1\end{pmatrix} + \kappa\begin{pmatrix}-C_{1,\beta} \\ -sq_2 \\ 0 \end{pmatrix} + \kappa^2\begin{pmatrix}\gamma_2 \\ \delta_2 \\ -\frac{1}{2}\left(C^2_{1,\beta} + sq^2_2\right)\end{pmatrix} + \\ \kappa^3\begin{pmatrix}
	-\mu \frac{b}{a+c}C_{1,\alpha} - \frac{b}{L}C_{1,\alpha}C_{2,\beta} - \frac{1}{L}C_{1,\beta}C_{2,h}\\
	\frac{a}{L}C_{1,\alpha}C_{2,\beta} + \frac{1}{L}C_{1,\alpha}C_{2,h} + \mu C_{1,\alpha} \\
	C_{1,\beta}\gamma_2 + sq_2\delta_2
	\end{pmatrix}
\end{aligned}
\end{equation}

Для третьего шатуна:

\begin{equation}
\begin{aligned}
	\frac{1}{L}S_3^0 = \begin{pmatrix} 0 \\ 0 \\ 1\end{pmatrix} + \kappa\begin{pmatrix}-C_{1,\beta} \\ -sq_3 \\ 0 \end{pmatrix} + \kappa^2\begin{pmatrix}\gamma_3 \\ \delta_3 \\ -\frac{1}{2}\left(C^2_{1,\beta} + sq^2_3\right)\end{pmatrix} + \\ \kappa^3\begin{pmatrix}
	-\mu \frac{b}{a+c}C_{1,\alpha} - \frac{1}{L}C_{1,\beta}C_{2,h}\\
	-\frac{c}{L}C_{1,\alpha}C_{2,\beta} + \frac{1}{L}C_{1,\alpha}C_{2,h} + \lambda C_{1,\alpha} \\
	C_{1,\beta}\gamma_3 + sq_3\delta_3
	\end{pmatrix}
\end{aligned}
\end{equation}

\section{Уравнения обратной задачи динамики}

Схема 3DOF платформы с изображением сил приведена на рис. 1. Особенностью используемой кинематической схемы является центральный карданный шарнир, удерживающий платформу от перемещений по курсу.

\begin{figure}[h]
	\center
	\includegraphics[width=0.75\textwidth]{3dof_kinematics_scheme_dynamic}
	\caption{Кинематическая схема подвижности 3DOF}
	\label{img:3dof_dynamic_scheme}
\end{figure}

Для решения обратной задачи динамики требуется определить силы реакции шатунов на платформу. Т.к. векторы данных сил всегда направлены вдоль шатунов, эти векторы для удобства можно представить в виде произведений $k_1S_1$, $k_2S_2$, $k_3S_3$, где $k_1$, $k_2$, $k_3$ – масштабные коэффициенты, $S_1$, $S_2$, $S_3$ – векторы положения шатунов в неподвижной системе координат, которые определяются по координатам точек крепления шатунов к кривошипам и подвижной платформе.

Обозначим за $S_1^0$, $S_2^0$, $S_3^0$ векторы от центра тяжести до точек крепления шатунов к подвижной платформе. Также ведем обозначения вертикальных проекций векторов $S_{iz}$, нормированных по длине шатуна в земной и подвижной системе координат:
\begin{equation}
\label{eq:DynamicVerticaldPojections}
	F_i = \frac{1}{L}S_{i,z};~ D_i = \frac{1}{L}S_{i,z}^0;~i=1\dots3
\end{equation}

Т.к. движение по курсу отсутствует, при составлении уравнения сил учитывается лишь вертикальная составляющая:
\begin{equation}
\label{eq:DynamicVerticalForcesEq1}
	m\ddot{h} = -k_sh + f_0 - mg + \sum_{i=1}^3 k_iS_{i,z},
\end{equation}
где $k_s$ -- коэффициент упругости пружины, $h$ -- вертикальное смещение, $m$ -- масса подвижной платформы с закрепленным на ней объектом, $f_0$ -- значение подъемной силы, создаваемой пружиной при $h=0$.

Перепишем уравнение относительно нормированных величин $F_i$:
\begin{equation}
\label{eq:DynamicVerticalForcesEq2}
	\sum_{i=1}^3k_iF_i = \frac{1}{L}\left(m\ddot{h}+k_sh-f_0+mg\right)
\end{equation}

Далее запишем уравнение моментов движения твердого тела в связанной системе координат. Уравнение моментов записывается относительно центра тяжести подвижной платформы с закрепленным объектом. Положение центра тяжести в системе координат платформы обозначим $X^0_C = \left[x_c, y_c, z_c\right]^T$. Тогда момент, создаваемый центральной опорой будет равен:
\begin{equation}
\label{eq:TorqueFromCentralSupport}
M_S^0 = -X_C^0 \times F_S^0 = -\left(X_C^0 \times e_3\right)\left(f_0 - k_sh\right),
\end{equation}

где $F^0_S$ – сила, действующая на платформу со стороны центральной опоры, $e_3 = \left[0~0~1\right]^T$ -- единичный вектор, направленный по оси $Z$. В дальнейшем будут использоваться первые две компоненты векторов моментов, т.к. вращение по курсу не рассматривается. Для того, чтобы учесть этот факт, введем вспомогательную матрицу $E_{12}$:
\begin{equation}
\label{eq:AuxiliaryUnitMatrix}
E_{12} = \begin{pmatrix}
1 & 0 & 0 \\
0 & 1 & 0
\end{pmatrix},
\end{equation}

Рассчитаем моменты в плоскости XY от каждого шатуна в системе координат платформы:
\begin{equation}
\label{eq:TorquesFromRodsInMovingCS}
\begin{aligned}
& E_{12}M_1^0 = E_{12}k_1\begin{pmatrix}
-a-x_c \\ b-y_c \\ -z_c
\end{pmatrix} \times S_1^0 = \\
& = k_1\left\{
\left[\begin{pmatrix}b \\ a\end{pmatrix} +
x_c\begin{pmatrix}0 \\ 1\end{pmatrix} -
y_c\begin{pmatrix}1 \\ 0\end{pmatrix}
\right]S_{1,z}^0+z_c\begin{pmatrix}S_{1,y}^0 \\ -S_{1,x}^0\end{pmatrix}\right\} \\
& E_{12}M_2^0 = E_{12}k_2\begin{pmatrix}
-a-x_c \\ -b-y_c \\ -z_c
\end{pmatrix} \times S_2^0 = \\
& = k_2\left\{
\left[\begin{pmatrix}-b \\ a\end{pmatrix} +
x_c\begin{pmatrix}0 \\ 1\end{pmatrix} -
y_c\begin{pmatrix}1 \\ 0\end{pmatrix}
\right]S_{2,z}^0+z_c\begin{pmatrix}S_{2,y}^0 \\ -S_{2,x}^0\end{pmatrix}\right\} \\
& E_{12}M_3^0 = E_{12}k_3\begin{pmatrix}
c-x_c \\ -y_c \\ -z_c
\end{pmatrix} \times S_3^0 = \\
& = k_3\left\{
\left[\begin{pmatrix}0 \\ c\end{pmatrix} +
x_c\begin{pmatrix}0 \\ 1\end{pmatrix} -
y_c\begin{pmatrix}1 \\ 0\end{pmatrix}
\right]S_{3,z}^0+z_c\begin{pmatrix}S_{3,y}^0 \\ -S_{3,x}^0\end{pmatrix}\right\} \\
\end{aligned}
\end{equation}

Запишем угловую скорость вращения платформы в подвижной системе координат. Положение произвольной точки $r$ в неподвижной СК определяется как $r^0 = R_xR_yr = Pr$. Тогда линейная скорость точек платформы $V = -\Omega\times\left(Pr\right) = \dot{P}r$. Равенство верно для любой точки $r$, поэтому справедлива замена $P^{-1}x$. Окончательно получим следующее выражение:

\begin{equation}
\label{eq:AngularSpeedFromLinear}
\Omega \times x = -\dot{P} P^{-1} x ~~\forall x\in\Re^3
\end{equation}

Раскроем правую часть равенства \ref{eq:AngularSpeedFromLinear}:

\begin{equation}
\label{eq:AngularSpeedTurnOperatorDiff}
\begin{aligned}
& \dot{P}P^{-1} = \left(\dot{R_x}R_y\dot{\alpha} + R_x\dot{R_y}\dot{\beta}\right)R_y^{-1}R_x^{-1}=\dot{\alpha}\dot{R_x}R_x^{-1} + \dot{\beta}R_x\dot{R_y}R_y^{-1}R_x^{-1}\\
& \dot{R_x} = \begin{pmatrix}
0 & 0 & 0\\
0 & -\sin{\alpha} & \cos{\alpha}\\
0 & -\cos{\alpha} & -\sin{\alpha}
\end{pmatrix},\\
& \dot{R_x}R_x^{-1} = \begin{pmatrix}
0 & 0 & 0\\
0 & 0 & 1\\
0 & -1 & 0
\end{pmatrix},~
 \dot{R_y}R_y^{-1} = \begin{pmatrix}
0 & 0 & -1\\
0 & 0 & 0\\
1 & 0 & 0
\end{pmatrix}.
\end{aligned}
\end{equation}

Окончательно, равенство \ref{eq:AngularSpeedTurnOperatorDiff} преобразуется к виду:

\begin{equation}
\label{eq:AngularSpeedTurnOperatorDiff2}
\Omega x = \begin{pmatrix}
0 & -\Omega_z & \Omega_y \\
\Omega_z & 0 & -\Omega_x \\
-\Omega_y & \Omega_x & 0
\end{pmatrix}x = \begin{pmatrix}
0 & -\dot{\beta}\sin{\alpha} & -\dot{\beta}\cos{\alpha} \\
\dot{\beta}\sin{\alpha} & 0 & \dot{\alpha} \\
\dot{\beta}\cos{\alpha} & -\dot{\alpha} & 0
\end{pmatrix}x
\end{equation}

С учетом матрицы $E_{12}$ вектор угловых скоростей имеет вид:

\begin{equation}
\label{eq:AngularSpeedInMovingCS}
\Omega^0 = \begin{pmatrix}
\Omega_x^0 \\ \Omega_y^0 \\ \Omega_z^0
\end{pmatrix} =
\begin{pmatrix}
\dot{\alpha}\cos{\beta} \\ \dot{\beta} \\ -\dot{\alpha}\sin{\beta}
\end{pmatrix},\qquad
E_{12}\Omega^0 =
\begin{pmatrix}
\dot{\alpha}\cos{\beta} \\ \dot{\beta}
\end{pmatrix}
\end{equation}

Уравнение моментов с учетом введенной в \ref{eq:AuxiliaryUnitMatrix} матрицы $E_{12}$ имеет вид:
\begin{equation}
\label{eq:TorquesEqCommon}
E_{12}\left[I\dot{\Omega}^0 + \Omega^0 \times \left[I\Omega^0\right] \right] =
E_{12}\left[M_1^0 + M_2^0 + M_3^0 + M_S^0\right]
\end{equation}

где $I$ – матрица моментов инерции платформы. Т.к. момент инерции записывается относительно центра тяжести, преобладающими являются диагональные компоненты:
\begin{equation}
I = \begin{pmatrix}
I_x & 0 & 0 \\ 0 & I_y & 0 \\ 0 & 0 & I_z
\end{pmatrix}
\end{equation}
Раскроем левую инерционную часть уравнения моментов \ref{eq:TorquesEqCommon}:
\begin{equation}
\label{eq:TorquesEqLeftInertialPart}
E_{12}\left[I\dot{\Omega}^0 + \Omega^0 \times \left[I\Omega^0\right] \right] = \begin{pmatrix}I_x\ddot{\alpha}\cos{\beta} \\ I_y\ddot{\beta}\end{pmatrix} +
\begin{pmatrix}\left(I_y-I_x-I_z\right)\dot{\beta} \\ \left(I_z-I_x\right)\dot{\alpha}\cos{\beta}\end{pmatrix}\dot{\alpha}\sin{\beta}
\end{equation}

Обозначим инерционные моменты $m_{\alpha}$ и $m_\beta$ для сокращения дальнейших выкладок:
\begin{equation}
\label{eq:TorquesInertial}
\begin{aligned}
& m_{\alpha} = I_x\ddot{\alpha}\cos{\beta} + \left(I_y-I_x-I_z\right)\dot{\alpha}\dot{\beta}\sin{\beta} \\
& m_{\beta} = I_y\ddot{\beta} + \left(I_z-I_x\right)\dot{\alpha}^2\sin{\beta}\cos{\beta}
\end{aligned}
\end{equation}

После подстановки \ref{eq:TorqueFromCentralSupport}, \ref{eq:TorquesFromRodsInMovingCS} и \ref{eq:TorquesInertial} в \ref{eq:TorquesEqCommon}, совместно с \ref{eq:DynamicVerticalForcesEq2} получим систему алгебраических уравнений относительно неизвестных коэффициентов $k_1$, $k_2$, $k_3$, которые в данном случае имеют смысл модулей векторов сил, действующих на шатуны:
\begin{equation}
\label{eq:KoefsLinearSystem1}
\left\{
\begin{aligned}
& F_1k_1 + F_2k_2 + F_3k_3 = \frac{1}{L}\left(m\ddot{h}+k_sh-f_0+mg\right), \\
& b\left(D_1k_1-D_2k_2\right)-Dky_c + S_y^0kz_c = \frac{1}{L}\left(m_{\alpha} + y_c\left(f_0 - k_sh\right)\right), \\
& a\left(D_1k_1+D_2k_2\right)-cD_3k_3+Dkx_c-S_x^0kz_c = \frac{1}{L}\left(m_{\beta} - x_c\left(f_0 - k_sh\right)\right),
\end{aligned}
\right.
\end{equation}

где $D=\left[D_1~D_2~D_3\right],~F=\left[F_1~F_2~F_3\right], k=\left[k_1~k_2~k_3\right]^T,~S^0_x=\left(S^0_{1,x},~S^0_{2,x},~S^0_{3,x}\right)/L, S^0_y=\left(S^0_{1,y},~S^0_{2,y},~S^0_{3,y}\right)/L$.

Прибавим ко второму уравнению в \ref{eq:KoefsLinearSystem1} правую часть первого, умноженную на $y_c$, к третьему –-- ту же правую часть, умноженную на $-x_c$. При этом для левой части второго и третьего уравнений заменим данное слагаемое на $Fk$ согласно равенству в первом уравнении. Таким образом преобразованные второе и третье уравнения системы \ref{eq:KoefsLinearSystem1}:
\begin{equation}
\label{eq:KoefsLinearSystemSimplification}
\begin{aligned}
& b\left(D_1k_1-D_2k_2\right)-Dky_c + S_y^0kz_c + Fky_c = \frac{1}{L}\left(m_{\alpha} + y_c\left(f_0 - k_sh\right)\right) + \\ & \frac{1}{L}\left(m\ddot{h}+k_sh-f_0+mg\right)y_c, \\
& a\left(D_1k_1+D_2k_2\right)-cD_3k_3+Dkx_c-S_x^0kz_c - Fkx_c = \\ & = \frac{1}{L}\left(m_{\beta} - x_c\left(f_0 - k_sh\right)\right) - \frac{1}{L}\left(m\ddot{h}+k_sh-f_0+mg\right)x_c,
\end{aligned}
\end{equation}

В \ref{eq:KoefsLinearSystemSimplification} выделим $E=D-F$ –- вспомогательную переменную, характеризующую горизонтальные смещения центра тяжести.

Упростив уравнение \ref{eq:KoefsLinearSystemSimplification}, вместе с первым уравнением  \ref{eq:KoefsLinearSystem1} получим преобразованную систему линейных уравнений относительно $k$:

\begin{equation}
\label{eq:KoefsLinearSystem2}
\begin{aligned}
& \begin{pmatrix}
F_1 & F_2 & F_3 \\
bD_1 & -bD_2 & 0 \\
aD_1 & aD_2 & -cD_3
\end{pmatrix}k +
\begin{pmatrix}0 \\ S_y^0 \\ -S_x^0\end{pmatrix}kz_c +
\begin{pmatrix}0 \\ -y_c \\ x_c \end{pmatrix}Ek = \frac{1}{L}
\begin{pmatrix}
m\ddot{h} \\
m_{\alpha} \\
m_{\beta}
\end{pmatrix} + \\
& + \frac{1}{L}\begin{pmatrix}
k_sh - f_0 + mg \\
y_c(m\ddot{h} + mg) \\
-x_c(m\ddot{h} + mg)
\end{pmatrix}
\end{aligned}
\end{equation}

Далее расчет коэффициентов $k$ осуществляется численно методом простых итераций.

Окончательно, получим формулы для определения вращающих моментов, составив уравнения для каждого кривошипа:
\begin{equation}
\label{eq:TorquesFromCrankAngles}
Q_i  = \begin{pmatrix}
0 \\ R\sin{q_i} \\ -R\cos{q_i}
\end{pmatrix} \times
\left(k_iS_i\right)
\end{equation}

После подстановки $S_i$ в \ref{eq:TorquesFromCrankAngles} получим формулы для моментов относительно рассчитанных значений $k_i$ :
\begin{equation}
\label{eq:TorquesFromCrankAnglesDetailed}
\begin{aligned}
& Q_i = k_iG_i,~i=1,\dots,3 \\
& G_1 = R \left\{ L\sin{q_1}-R\sin{q_1}\cos{q_1}-R^2 \left[ \frac{1}{8b} {\left( \cos{q_1}-\cos{q_2} \right)}^2 \cos{q_1} + \frac{1}{2L}\sin^3{q_1} \right] \right. +\\
& \left.+ \frac{R^3}{8Lb} \left( \cos{q_1}-\cos{q_2} \right)^2 \left(\cos{2q_1}+\cos{q_1}\cos{q_2}\right)
\right\} \\
& G_2 = R \left\{ L\sin{q_2}-R\sin{q_2}\cos{q_2}+R^2 \left[ \frac{1}{8b} {\left( \cos{q_1}-\cos{q_2} \right)}^2 \cos{q_2} - \frac{1}{2L}\sin^3{q_2} \right] \right. -\\
& \left.- \frac{R^3}{8Lb} \left( \cos{q_1}-\cos{q_2} \right)^2 \left(\cos{2q_2}+\cos{q_1}\cos{q_2}\right)
\right\} \\
& G_3 = R\left\{L\sin{q_3}-R\sin{q_3}\cos{q_3}-R^2\frac{1}{2L}\sin^3q_3\right\}
\end{aligned}
\end{equation}

\section{Исследование решения линейной системы относительно неизвестных сил реакции шатунов}

\subsection{Обращение матрицы линейной системы}
Матрица системы \ref{eq:LinearSystemStatic}:
\begin{equation}
\begin{aligned}
	A_0 =
	\begin{pmatrix}
	F_1 & F_2 & F_3 \\
	bD_1 & -bD_2 & 0 \\
	aD_1 & aD_2 & -cD_3
	\end{pmatrix}
\end{aligned}
\end{equation}

обращается явно:
\begin{equation}
\begin{aligned}
	A_0^{-1} = \frac{1}{\Delta_0}\begin{pmatrix}
	bcD_2D_3 & cD_3F_2 + aD_2F_3 & bD_2F_3 \\
	bcD_1D_3 & -aD_1F_3-cD_3F_1 & bD_1F_3 \\
	2abD_1D_2 & a\left(D_1F_2-D_2F_1\right) & -b\left(D_1F_2 + D_2F_1\right)
\end{pmatrix}
\end{aligned}
\end{equation}

где определитель матрицы $A_0$ равен
\begin{equation}
	\Delta_0 = 2abD_1D_2F_3 + bcD_3\left(D_1F_2 + D_2F_1\right)
\end{equation}

Для решения системы применим метод простой итерации. Запишем уравнения в виде:

\begin{equation}
\begin{aligned}
	A_0k + Hk = B, ~ H = \begin{pmatrix}0 \\ S_y^0 \\ -S_x^0 \end{pmatrix}z_c + \begin{pmatrix}0 \\ -y_c \\ x_c \end{pmatrix}E
\end{aligned}
\end{equation}

Пусть спектральный радиус $\rho\left(A_0^{-1}H\right) < 1$. Тогда решение можно представить в виде ряда:

\begin{equation}
\begin{aligned}
	k & = \left(A_0 + H\right)^{-1}B = A_0^{-1}\left(I + HA_0^{-1}\right)^{-1}B \\
	&= A_0^{-1}\sum_{i=0}^{\inf}\left(-HA_0^{-1}\right)^i B = A_0^{-1}\left(\sum_{i=0}^{\inf}C^i\right)B,
\end{aligned}
\end{equation}

где

\begin{equation}
\begin{aligned}
	C = -H A_0^{-1} = \begin{pmatrix}0 \\ S_y^0 \\ -S_x^0 \end{pmatrix}A_0^{-1}z_c + \begin{pmatrix}0 \\ -y_c \\ x_c \end{pmatrix}EA_0^{-1}
\end{aligned}
\end{equation}

Второе слагаемое мало, т.к. величина $E$ пропорциональна $\kappa^2$, а смещение $x_c,~y_c$ составляют сотые доли метра. Первое слагаемое рассчитывается в следующем разделе.

Особый интерес представляет аппроксимация ряда частичной суммой первого или второго порядка:

\begin{equation}
\begin{aligned}
	\hat{k}^{(1)} = A_0^{-1}B + A_0^{-1}CB, ~~ \hat{k}^{(2)} = A_0^{-1}B + A_0^{-1}CB  + A_0^{-1}C^2B
\end{aligned}
\end{equation}

Запишем матрицу $C$ в виде:

\begin{equation}
\begin{aligned}
	C = -C_0z_c - \begin{pmatrix}0 \\ -y_c \\ x_c \end{pmatrix}EA_0^{-1}, ~C_0 = \begin{pmatrix}0 \\ S_y^0 \\ -S_x^0 \end{pmatrix}A_0^{-1}
\end{aligned}
\end{equation}

Назовем $C$ матрицей итераций.

\subsection{Расчет матрицы итераций}

Введем обозначение для нормированного определителя матрицы $A_0$:
\begin{equation}
\begin{aligned}
	\Delta = \Delta_0/b = 2aD_1D_2F_3 + cD_3\left(D_1F_2 + D_2F_1\right)
\end{aligned}
\end{equation}

и для столбцов нормированной присоединенной матрицы $A_1 = \Delta A_0^{-1}e_1,~A_2 = \Delta A_0^{-1}e_2,~A_3 = \Delta A_0^{-1}e_3$:

\begin{equation}
\begin{aligned}
	A_1 = \begin{pmatrix}cD_2D_3 \\ cD_1D_3 \\ 2aD_1D_2\end{pmatrix},~ A_2 = \frac{1}{b}\begin{pmatrix}cD_3F_2 + aD_2F_3 \\ -aD_1F_3-cD_3F_1 \\ a\left(D_1F_2-D_2F_1\right)\end{pmatrix}, ~A_3=\begin{pmatrix}D_2F_3 \\ D_1F_3 \\ -\left(D_1F_2 + D_2F_1\right) \end{pmatrix}
\end{aligned}
\end{equation}

Вычислим каждый столбец в матрице:

\begin{equation}
\begin{aligned}
	C_0 = \begin{pmatrix}0 \\ S_y^0 \\ -S_x^0\end{pmatrix}A_0^{-1} = \frac{1}{\Delta}\begin{pmatrix}
	0 & 0 & 0 \\ S_y^0A_1 & S_y^0A_2 & S_y^0A_3 \\ -S_x^0A_1 & -S_x^0A_2 & -S_x^0A_3
	\end{pmatrix}
\end{aligned}
\end{equation}

С точностью до третьего порядка:

\begin{equation*}
	\begin{aligned}
		\frac{1}{\Delta}\begin{pmatrix}S_y^0A_1 \\ -S_x^0A_1\end{pmatrix} & = \frac{2\left(a+c\right)}{\Delta}\left\lbrace\begin{pmatrix}\sin\alpha\cos\beta \\ \sin\beta\end{pmatrix}- \right.\\
		& \left. -\kappa\begin{pmatrix}1\\0\end{pmatrix}\frac{2a\sin{q_3}+c\sin{q_1}+c\sin{q_2}}{2\left(a+c\right)}\left(\cos\alpha-\kappa^2C_{1,\beta}^2\right) - \right.\\
		& \left. -\frac{1}{2}\kappa^2 \begin{pmatrix} \sin{\alpha} \\ \sin{\beta} \end{pmatrix} \frac{2a\sin^2{q_3}+c\left(\sin^2{q_1}+\sin^2{q_2}\right)}{2\left(a+c\right)}-\kappa^2 \begin{pmatrix} \sin{\alpha} \\ 0 \end{pmatrix} C^2_{1,\beta} + \right.\\
		& \left. \frac{1}{2} \kappa^3 \begin{pmatrix} 1 \\ 0 \end{pmatrix} \frac{2asq_3 \left(sq_1^2+sq_2^2 \right) + csq_1\left( sq_2^2+sq_3^2 \right) + csq_2\left(sq_1^2+sq_3^2\right)}{2\left(a+c\right)} \right\rbrace\\
		 S_y^0 A_2 & \approx -\kappa \frac{a+c}{b} \left( \sin{q_1} - \sin{q_2} \right) \left\lbrace 1 + \frac{1}{2} \kappa^2 \left(4\lambda - C^2_{1,\beta} + \sin{q_1} \sin{q_2} - \sin^2{q_3}\right) \right\rbrace -\\
		& -\kappa^2\frac{a+c}{L}C_{1,\alpha}\\
		S_y^0 A_3 & \approx \kappa\left(2\sin{q_3} - \sin{q_1} - \sin{q_2}\right) \left(1 - \frac{1}{2}\kappa^2\left(C^2_{1,\alpha} + C^2_{1,\beta}\right)\right) + \\
		& + \frac{1}{2}\kappa^3 \left\lbrace \sin{q_1}\left(\sin^2{q_3} + sq^2_2\right) + \sin{q_2}\left(\sin^2{q_3} + sq^2_1\right)\right.\\
		&\left. -\sin{q_3}\left(\sin^2{q_1} + \sin^2{q_2} + sq^2_1 + sq^2_2\right)\right\rbrace \\
		S^0_x A_2 & \approx \frac{1}{2b}\kappa^3 C_{1,\beta}c\left(\sin^2{q_2} - \sin^2{q_1}\right) \\
		S^0_x A_3 & \approx \frac{1}{2}\kappa^3 C_{1,\beta}\left(2\sin^2{q_3} - \sin^2{q_1} - \sin^2{q_2}\right) + \kappa^2\frac{a+c}{L}C^2_{1,\beta} + \\
		& + \kappa^2\frac{a+c}{L}\left(C^2_{1,\beta} + 2\kappa C_{1,\beta}C_{2,\beta}\right)
	\end{aligned}
\end{equation*}

Т.к. $\kappa^3$ является малой величиной и умножается на линейную комбинацию тригонометрических функций, компонентами матрицы итераций, содержащими $\kappa^3$ можно пренебречь. Запишем матрицу итераций до второго порядка $\kappa$:

\begin{equation}
\begin{aligned}
	C_0 & \approx \frac{2\left(a+c\right)}{\Delta}\left\lbrace \begin{pmatrix}
	0 & 0 & 0 \\ \sin{\alpha}\cos{\beta} & 0 & 0 \\ \sin{\beta} & 0 & 0
	\end{pmatrix} + \right.\\
	&\left. + \kappa\begin{pmatrix}
	0 & 0 & 0 \\
	-\frac{2a\sin{q_3}+c\sin{q_1}+c\sin{q_2}}{2\left(a+c\right)} & \frac{\sin{q_2}-\sin{q_1}}{2b} & \frac{2\sin{q_3} - \sin{q_1} - \sin{q_2}}{2\left(a+c\right)}\\
	0 & 0 & 0 \\
	\end{pmatrix} - \right.\\
	& \left.\frac{1}{2L}\kappa^2\begin{pmatrix}
	0 & 0 & 0 \\
	0 & C_{1,\alpha} & 0 \\
	0 & 0 & C^2_{1,\beta}
	\end{pmatrix}\right\rbrace
\end{aligned}
\end{equation}
 Определитель $\Delta$ отличается от $2\left(a+c\right)$ на величину порядка $\kappa^2$. При возведении в квадрат матрицы $z_cC_0$ умножаются на 0 все элементы первого столбца. Множитель $\kappa^2$ останется только в центральном элементе:

\begin{equation}
\begin{aligned}
	\kappa^2\left(\frac{2\left(a+c\right)}{\Delta}\right)^2 \frac{z^2_c}{\left(2b\right)^2}\left(\sin{q_2} - \sin{q_1}\right)
\end{aligned}
\end{equation}

которым, как правило можно пренебречь.

Поэтому одна итерация в решении системы линейных уравнений дает неплохую аппроксимацию решения.

\subsection{Слагаемое с горизонтальным смещением центра тяжести}

Второе слагаемое в матрице итераций $C$ содержит произведение вектора из смещений $x_cy_c$ на коэффициент $EA_0^{-1}$, где $E = D - F$.

Оценим компоненты строки $E = \left(E_1,~E_2,~E_3\right)$. При $i=1,2,3$:

\begin{equation}
\begin{aligned}
	E_i & \approx \frac{1}{2}\kappa^2 \left(-q^2_i - C^2_{1,\beta} + \sin^2{q_i}\right) + \kappa^3 \left\lbrace C_{1,\beta}\gamma_i + sq_i\delta_i + \frac{b_i}{2L}\sin{q_i}C^2_{1,\alpha} \right\rbrace \\
	& + \kappa^2\left[C_{1,\alpha}\sin{q_i} - \frac{1}{2}\left(C^2_{1,\alpha} + C^2_{1,\beta}\right)\right] + \\
	& + \kappa^3\left\lbrace C_{2,\alpha}\sin{q_i} + \frac{a}{2L}C^3_{1,\beta} - \left(C_{1,\beta}C_{2,\beta} + C_{1,\beta}C_{2,\alpha}\right)\right\rbrace
\end{aligned}
\end{equation},

где $b_i$ обозначает следующее: $b_1 = b,~b_2 = -b,~b_3 = 0$. Компоненты строки $E$ имеют второй порядок относительно величины $\kappa$.

Оценим каждую компоненту коэффициента $EA^{-1}_0$. Первая компонента:
\begin{equation}
\begin{aligned}
EA^{-1}_0e_1 \approx & \kappa^2\left[\left(C_{1,\alpha} + \kappa C_{2,\alpha} - \kappa^2 C^2_{1,\beta}\right)\frac{2a\sin{q_3} + c\left(\sin{q_1} + \sin{q_2}\right)}{\Delta} - \right.\\
&\left. - \frac{a+c}{\Delta}\left(C^2_{1,\alpha} + C^2_{1,\beta}\right)\left(1 - \kappa^2 C^2_{1,\beta}\right)\right] + \\
& + \kappa^3\frac{2\left(a+c\right)}{\Delta} \left[\frac{a}{2L}C^3_{1,\beta} - \left(C_{1,\beta}C_{2,\beta} + C_{1,\alpha}C_{2,\alpha}\right)\right] + \\
& + \frac{1}{2}\kappa^4\left[ \frac{1}{2}\left(C^2_{1,\alpha} + C^2_{1,\beta}\right) \frac{2a\left(sq^2_1 + sq^2_2\right) + c\left(q^2_1 + q^2_2 + 2q^2_3\right)}{\Delta}\right. \\
& -\left. \frac{2a\sin{q_3}\left(q^2_1 + q^2_2\right) + c\sin{q_1}\left(q^2_2 + q^2_3\right) + c\sin{q_2}\left(q^2_1 + q^2_3\right)}{\Delta} \right]
\end{aligned}
\end{equation}

Вторая компонента:
\begin{equation}
\begin{aligned}
EA^{-1}_0 e_2 & = \frac{1}{\Delta}EA_2 = \frac{a+c}{\Delta b}D_3\left(D_1F_2 - D_2F_1\right) \approx\\
& \approx \frac{2\left(a+c\right)}{\Delta}\kappa^2D_3C_{1,\alpha}\left\lbrace \frac{\sin{q_1} - \sin{q_2}}{2b}\left(1 - \frac{1}{2}\kappa C_{2,\alpha}\right) + \kappa \frac{1}{4L}\left(2C^2_{1,\beta} + C^2_{1,\alpha}\right)\right\rbrace.
\end{aligned}
\end{equation}

Третья компонента:
\begin{equation}
\begin{aligned}
EA^{-1}_0 e_3 & \approx -\kappa\sin{\alpha}\cos^2{\beta}\frac{2\sin{q_3} - \sin{q_1} - \sin{q_2}}{\Delta} - \\
& -\frac{1}{2\Delta}\kappa^4 \left\lbrace -\frac{1}{2}C^2_{1,\beta}\left(2\sin^2{q_3} - \sin^2{q_1} - \sin^2{q_2}\right) \right. \\
& \left. -C^2_{1,\alpha}\left(\sin{q_3} - \sin{q_1}\right)\left(\sin{q_3} - \sin{q_2}\right)\right. \\
& \left. +C_{1,\alpha}\left(\sin{q_1} - \sin{q_3}\right)\left(sq^2_2 - \sin{q_1}\sin{q_3}\right)\right. \\
& \left. -C_{1,\alpha}\left(\sin{q_2} - \sin{q_3}\right)\left(sq^2_1 - \sin{q_2}\sin{q_3}\right)\right\rbrace
\end{aligned}
\end{equation}


\section{Оценка параметров уравнений обратной задачи динамики}
Практически все параметры в \ref{eq:KoefsLinearSystem2} определить на практике с достаточной точностью не представляется возможным из-за сложностей в организации процесса измерения. В связи с этим, целесообразно произвести параметрическую идентификацию, используя данные с датчиков, установленных на СП.

\subsection{Идентификация параметров в статике}
Параметрами, в большей степени влияющими на результат расчета движущих моментов являются вес закрепленного на платформе объекта, коэффициент жесткости $k_s$, начальное натяжение пружины $f_0$ и смещение центра тяжести $(x_c, y_c, z_c)$. Эти величины не меняются в течение длительного времени, поэтому целесообразно произвести идентификацию объекта на основе данных эксперимента.
	Данные параметры можно определить в статическом режиме, т.е. измерив вращающие моменты в нескольких стационарных положениях платформы.

В стационарном положении $\dot{\alpha} = \dot{\beta} = \ddot{\alpha} = \ddot{\beta} = \dot{h} = \ddot{h} = 0$, поэтому в правой части \ref{eq:KoefsLinearSystem2} обращаются в нуль все три инерционных слагаемых:
\begin{equation}
\label{eq:StaticLimit}
m\ddot{h} \to 0, \qquad
m_{\alpha} = I_x\ddot{\alpha}\cos{\beta} + \left(I_y-I_x-I_z\right)\dot{\alpha}\dot{\beta}\sin{\beta} \to 0, \qquad
m_{\beta} = I_y\ddot{\beta} + \left(I_z-I_x\right)\dot{\alpha}^2\sin{\beta}\cos{\beta} \to 0,
\end{equation}
а слагаемые $y_c(m\ddot{h} + mg)$ и $-x_c(m\ddot{h} + mg)$ переходят в $y_cmg$ и $-x_cmg$. Левая часть при этом не меняется: две матрицы при $z_c$ и при $E$ лишь объединяются в одну,
\begin{equation}
\label{eq:StaticRegroup}
\begin{pmatrix}0 \\ S_y^0 \\ -S_x^0\end{pmatrix}kz_c +
\begin{pmatrix}0 \\ -y_c \\ x_c \end{pmatrix}Ek =
\begin{pmatrix}0 \\ S_y^0z_c - y_cE \\ -S_x^0z_c + x_cE\end{pmatrix}k .
\end{equation}

Таким образом, уравнения системы \ref{eq:KoefsLinearSystem2} путем обнуления производных принимают вид:
\begin{equation}
\label{eq:LinearSystemStatic}
\begin{aligned}
& \begin{pmatrix}
\bar{F_1} & \bar{F_2} & \bar{F_3} \\
b\bar{D_1} & -b\bar{D_2} & 0 \\
a\bar{D_1} & a\bar{D_2} & -c\bar{D_3}
\end{pmatrix}\bar{k} +
\begin{pmatrix}0 \\ \bar{S}_y^0z_c -y_c\bar{E} \\ -\bar{S}_x^0z_c+x_c\bar{E}\end{pmatrix}\bar{k} = \frac{1}{L}
\begin{pmatrix}
k_s\bar{h} - f_0 + mg \\
y_cmg \\-x_cmg
\end{pmatrix}
\end{aligned}
\end{equation}

Переменные с чертой в \ref{eq:LinearSystemStatic} обозначают константы, рассчитанные по данным эксперимента или полученные из исходных данных, без черты – идентифицируемые параметры \cite{Graupe}.

Введем дополнительные расчетные величины, которые в \ref{eq:LinearSystemStatic} не зависят от идентифицируемых параметров:
\begin{equation}
\label{eq:AdditionalEstimatedVariables}
\begin{pmatrix}
\bar{p_1} \\ \bar{p_2} \\ \bar{p_3}
\end{pmatrix} =
\begin{pmatrix}
\bar{F_1} & \bar{F_2} & \bar{F_3} \\
b\bar{D_1} & -b\bar{D_2} & 0 \\
a\bar{D_1} & a\bar{D_2} & -c\bar{D_3}
\end{pmatrix}\bar{k}
\end{equation}

Тогда после подстановки \ref{eq:AdditionalEstimatedVariables} в \ref{eq:LinearSystemStatic} получим:
\begin{equation}
\label{eq:AdditionalParamsRightPart}
\begin{aligned}
& \bar{p_1} = \frac{1}{L}\left(mg - f_0 + k_sh\right)\\
& \bar{p_2} = \left(\frac{mg}{L}+\bar{E}\bar{k}\right)y_c - \bar{S}^0_y\bar{k}z_c \\
& \bar{p_3} = \bar{S}^0_x\bar{k}z_c - \left(\frac{mg}{L}+\bar{E}\bar{k}\right)x_c
\end{aligned}
\end{equation}

Пусть по результатам эксперимента были получены n стационарных положений платформы. Т.к. в процессе эксперимента вращающие моменты $Q_1,~Q_2,~Q_3$ измеряются, а $G_1,~G_2,~G_3$ рассчитываются по кинематическим уравнениям, вычисляется вектор $k$ путем решения системы (11) методом простых итераций. Далее на каждом из участков по формуле (18) вычисляются значения  $p_{1,i},~p_{2,i},~p_{3,i},~i = 1,\dots,n$.

Вместо величин $\left(x_c, y_c\right)$ определяются сразу произведения $\left(x_c, y_c\right) = \left(x_c, y_c\right) mg/L$.

Для предварительной (раздельной) оценки параметров применяется метод наименьших квадратов с весовыми множителями \cite{Graupe}. Невязки записываются непосредственно по \ref{eq:AdditionalParamsRightPart}:

\begin{equation}
\label{eq:LeastSquareMethodFormulation}
\begin{aligned}
& J_{xyz}\left(x_c,y_c,z_c\right)=\sum_{i=1}^n\left[
c_{i,x}\left|\bar{p}_{3,i}-\left(\bar{S}_x^0\bar{k}\right)_iz_c + \bar{x_c}\right|^2
+ c_{i,y}\left|\bar{p}_{2,i}+\left(\bar{S}_y^0\bar{k}\right)_iz_c - \bar{y_c}\right|^2\right]
\rightarrow\min, \\
& J_{sf}\left(k_s,~f_0\right)=\sum_{j=1}^m{\left|\bar{p}_{1,j}-\frac{k_s\bar{h_j}}{L}-\frac{mg-f_0}{L}\right|}^2
\rightarrow\min.
\end{aligned}
\end{equation}

где $c_{i,x}$ и $c_{i,y}$ -- весовые множители, $i = 1,\dots,n$ -- номера стационарных участков, вошедших в оценивание смещений центра тяжести, $j = 1,\dots,m$ -- номера участков, вошедших в оценивание параметров пружины.

Обратим внимание на множитель $1/L$ при $k_s$. Уравнение сил \ref{eq:DynamicVerticalForcesEq2} даёт $\bar{p}_1 = \left(k_s\bar{h} + mg - f_0\right)/L$, поэтому коэффициентом при $k_s$ служит $\bar{h}/L$, а не $\bar{h}$. Если это деление опустить, оценивается величина $k_s/L$, а не сама жёсткость пружины.

Раздельная схема требует, чтобы вектор $\bar{k}$ вычислялся непосредственно из измерений, $\bar{k}_i = \bar{m}_i/\bar{G}_i$, и потому неприменима в конфигурациях, где кривошип находится вблизи мёртвой точки и $G_i \to 0$. Основной (групповой) метод, свободный от этого ограничения, рассмотрен в разделе \ref{sec:GroupMethod}.

\subsection{Идентификация в режиме гармонических колебаний}
Предположим, что в записи имеются несколько участков, на которых углы крена и тангажа меняются периодически. Экспериментально установлено, что среди инерциальных моментов сил только чистые вторые производные от углов крена и тангажа могут принимать значения, влияющие на общий суммарный момент сил. Это значит, что:

\begin{equation}
\begin{aligned}
m_{s,\alpha} & \approx 0,\\
m_{s,\beta} & \approx 0,\\
\left(I_y - I_x - I_z\right)\dot{\alpha}\dot{\beta}\sin{\beta} & \approx 0,\\
\left(I_z - I_x\right)\dot{\alpha}^2 \sin{\beta}\cos{\beta} & \approx 0.
\end{aligned}
\end{equation}

Наиболее значимый вклад вносит компонента $\begin{pmatrix}0 & I_x\ddot{\alpha}\cos{\beta} & I_y \ddot{\beta} \end{pmatrix}^T$.

Поскольку смещения центра тяжести и параметры сил оценены по стационарным участкам, в каждый момент времени коэффициенты сил $k$ можно представить в виде:
\begin{equation}
\begin{aligned}
	k = k_0 + k_{\alpha} + k_{\beta}
\end{aligned}
\end{equation}

где
$k_0 = A^{-1}B,~k_{\alpha}=A^{-1}\frac{1}{L}\begin{pmatrix}0 \\ I_x\ddot{\alpha}\cos{\beta} \\ 0\end{pmatrix},~k_{\beta}=A^{-1}\frac{1}{L}\begin{pmatrix}0 \\ 0 \\ I_y\ddot{\beta}\end{pmatrix},~$

и

\begin{equation}
\begin{aligned}
	A & = \begin{pmatrix}F_1 & F_2 & F_3 \\ bD_1 & -bD_2 & 0 \\ aD_1 & aD_2 & -cD_3\end{pmatrix}+
	\begin{pmatrix}0 \\ S_y^0 \\ -S_x^0\end{pmatrix}z_c + \begin{pmatrix}0 \\ -y_c \\ x_c\end{pmatrix}E, \\
	B &= \begin{pmatrix}
	\frac{1}{L}m\ddot{h} + \frac{1}{L}k_sh + \bar{f}\\
	\bar{y_c}\left(\ddot{h}/g + 1\right) \\
	-\bar{x_c}\left(\ddot{h}/g + 1\right)
	\end{pmatrix}
\end{aligned}
\end{equation}

Здесь $\bar{f} = \left(mg - f_0\right)/L$, $\bar{x_c} = x_c\,mg/L$, $\bar{y_c} = y_c\,mg/L$ -- константы, оценённые на стационарных участках.

Оценка моментов сил на валах кривошипов раскладывается на три компоненты:

\begin{equation}
\begin{aligned}
	m = m_0 + m_{\alpha} + m_{\beta},~m_0 = Gk_0,~m_{\alpha} = Gk_{\alpha},~m_{\beta} = Gk_{\beta}
\end{aligned}
\end{equation}

На каждом однородном участке колебаний с постоянной частотой это простые гармоники. При помощи БПФ определяется частота $F$ простых гармоник вращающих моментов. Величина $F$ измеряется в количестве периодов на длину выбранного участка сигнала $N$.
Далее определяются комплексные амплитуды этих гармоник:

\begin{equation}
\begin{aligned}
	m_0(t) = \operatorname{Re}\left(A_0 \exp{\frac{2\pi i}{N}Ft}\right),\\
	m_{\alpha}(t) = \operatorname{Re}\left(A_{\alpha} \exp{\frac{2\pi i}{N}Ft}\right),\\
	m_{\beta}(t) = \operatorname{Re}\left(A_{\beta} \exp{\frac{2\pi i}{N}Ft}\right).
\end{aligned}
\end{equation}

Каждая амплитуда $A_0, A_{\alpha}, A_{\beta}$ имеет 3 компоненты, соответствующие кривошипам. Аналогично выделяется первая гармоника измеряемой функции моментов на валах $m_{obs}$. Вектор её амплитуд обозначим $A_{obs}$. Неточно известны множители, с которыми линейные и угловые ускорения влияют на моменты на валах. Это масса кабины m и её моменты инерции $I_x$ и $I_y$. Однако влияние линейных ускорений $h$ на порядок меньше, чем угловых ускорений. Поэтому вместо $m$ можно подставить приблизительную оценку этой величины. Рассмотрим задачу настройки параметров $I_x$ и $I_y$. В измерениях $m_{obs}(t)$ возможны запаздывания, которые проявляются в изменении фазы в комплексной амплитуде $A_{obs}$. Поскольку смещение фазы неизвестно, то настройку будем проводить только по вещественным амплитудам гармоник. Пусть $i = 1,\ldots,n$ — номера участков сигнала с гармоническим изменением углов кривошипов. Определим функционал качества:

\begin{equation}
\begin{aligned}
	J\left(c_x,c_y\right) = \sum^n_{i=1}\left(\|A_{obs,i}\|^2 - \|A_{0,i} + c_xA_{\alpha,i} + c_yA_{\beta,i}\|^2\right)^2
\end{aligned}
\end{equation}

Этот функционал минимизируется по $\left(c_x, c_y\right)$ градиентным методом. При расчете предварительно задаются границы $c_{\min},~c_{\max}$ допустимых отклонений коэффициентов от 1.

После того, как коэффициенты настроены, получаем оценки комплексных амплитуд гармоник моментов на валах:

\begin{equation}
\begin{aligned}
	\hat{A}_i = A_{0,i} + c_xA_{\alpha,i} + c_yA_{\beta,i},~i=1,\ldots,n
\end{aligned}
\end{equation}

Для каждого $i$ это трехмерный вектор с компонентами, относящимися к каждому кривошипу.
Результаты идентификации, приведенные в главе 6, показывают, что измерения имеют существенное запаздывание по фазе, зависящие от частоты.
%
%По результатам идентификации для платформы с исходными данными $a = 0,25 м; b = 0,4 м; c = 0,5 м; R = 0,15 м; L = 0,56 м$ были определены следующие параметры модели:
%\begin{equation}
%\begin{aligned}
%	....
%\end{aligned}
%\end{equation}

\section{Групповой метод оценивания}
\label{sec:GroupMethod}

Раздельная схема предыдущего раздела требует вычислять $\bar{k}_i = \bar{m}_i/\bar{G}_i$ и потому неустойчива при $G_i \to 0$. Групповой метод снимает это ограничение: усилия в шатунах переводятся из вычисляемых величин в оцениваемые, а связь с измеренными моментами накладывается как мягкое ограничение.

Введём вектор неизвестных
\begin{equation}
\tau = \operatorname{col}\left(k_1,\ldots,k_n,~k_s,~\bar{f},~\bar{x_c},~\bar{y_c}\right)
\end{equation}
размерности $3n+4$ и будем минимизировать
\begin{equation}
\label{eq:GroupFunctional}
J\left(\tau, z_c\right) =
\underbrace{\sum_{i=1}^{n}\left\|A_i\left(z_c\right)k_i + M_i\theta\right\|^2}_{\text{блок 1: равновесие}}
+ \gamma^2\underbrace{\sum_{i=1}^{n}\left\|m_i - G_i\circ k_i\right\|^2}_{\text{блок 2: измерения}},
\end{equation}
где $\theta = \operatorname{col}\left(k_s,\bar{f},\bar{x_c},\bar{y_c}\right)$, $\circ$ -- покомпонентное умножение,
\begin{equation}
A_i\left(z_c\right) = A_{0,i} + \begin{pmatrix}0 \\ S^0_{y,i} \\ -S^0_{x,i}\end{pmatrix}z_c,
\qquad
M_i = -\begin{pmatrix}
\bar{h}_i/L & 1 & 0 & 0\\
0 & 0 & 0 & 1\\
0 & 0 & -1 & 0
\end{pmatrix}.
\end{equation}
При фиксированном $z_c$ задача линейно-квадратичная и решается явно; $z_c$ перебирается по сетке.

Существенно, что \textbf{блок 1 сам по себе не содержит информации о параметрах}. Каждая матрица $A_i$ обратима (для исследуемого стенда $\operatorname{cond}A_{0} \approx 3$ во всех рабочих конфигурациях), поэтому при любом $\theta$ можно положить $k_i = -A_i^{-1}M_i\theta$ и обратить блок 1 в нуль. Вся информация о $\theta$ приходит из блока 2, а множитель $\gamma$ переключает между двумя оценивателями:
\begin{itemize}
\item $\gamma \to \infty$: $k_i = m_i/G_i$ берётся из измерений, $\theta$ подгоняется под уравнения равновесия (equation error). Это в точности раздельная схема \ref{eq:LeastSquareMethodFormulation}; невязка по моментам тождественно равна нулю.
\item $\gamma \to 0$: $k_i$ следует уравнениям равновесия, $\theta$ подгоняется так, чтобы расчётные моменты $G_i\circ k_i$ совпали с измеренными (output error). Невязка по моментам остаётся и поддаётся интерпретации.
\end{itemize}
На идеальной модели оба предела дают одно и то же. Расхождение между ними служит мерой неадекватности модели и потому подлежит публикации наряду с самими оценками.

\section{Критерий отбора стационарных участков}

Момент на валу $i$-го кривошипа равен $Q_i = k_iG_i$, причём $G_i \propto \sin{q_i}$. При $q_i \to 0$ кривошип приходит в мёртвую точку, $G_i \to 0$, и момент на валу обращается в нуль независимо от усилия в шатуне. Такие конфигурации не несут информации об оцениваемых параметрах и должны исключаться.

Вместо перечисления номеров участков вручную используется формальный критерий: участок допускается к оцениванию, если
\begin{equation}
\label{eq:SegmentCriterion}
\min_{i=1,2,3}\left|G_i\right| > G_{\min}.
\end{equation}
Для исследуемого стенда при $G_{\min} = 0{,}02$ критерий отбирает те же участки $11-24$, что были выбраны ранее по содержательным соображениям. На отобранных участках $\left|G_i\right| = 0{,}05\ldots0{,}08$, на отброшенных -- до $10^{-5}$. Средняя ошибка воспроизведения моментов составляет $23$~Н$\cdot$м на отобранных участках против $75$~Н$\cdot$м на отброшенных при максимуме $287$~Н$\cdot$м.

Отдельно отметим участок $11$ -- первый после выхода из мёртвой точки. Невязки на нём ($+59{,}8$, $-66{,}4$, $+29{,}6$~Н$\cdot$м) втрое превышают типичные, что согласуется с наблюдением о непопадании в общую регрессионную зависимость при первом переходе от $q = 0$.

\section{Потери в приводе}

\subsection{Место включения в модель}

Величина $Q_i = k_iG_i$ есть момент, который механизм требует от вала кривошипа. Измеряется же электромагнитный момент двигателя, пересчитанный через передаточное число редуктора. Потери в редукторе и собственные механические потери двигателя лежат между этими двумя величинами, поэтому модель измерения принимает вид
\begin{equation}
\label{eq:FrictionModel}
Q_i^{\text{изм}} = k_iG_i + M_{\text{тр},i}\operatorname{sign}\dot{q}_i + J_{\text{прив}}\ddot{q}_i ,
\end{equation}
где $M_{\text{тр},i}$ -- момент потерь, $J_{\text{прив}}$ -- приведённая к валу кривошипа инерция ротора (см. следующий раздел). Слагаемые ортогональны: первое синфазно скорости, второе -- ускорению.

\subsection{Проверка на гармонических участках}

Невязка $r = Q^{\text{изм}} - kG$ проверена против трёх конкурирующих моделей, каждая с двумя параметрами:
\begin{equation}
r = c_0 + F_c\operatorname{sign}\dot{q}, \qquad
r = c_0 + C_v\dot{q}, \qquad
r = c_0 - \tau\,\frac{d}{dt}\left(kG\right).
\end{equation}
Решающим критерием служит не $R^2$ на отдельном участке, а постоянство параметра между частотами: истинный механизм даёт одно и то же значение на всех частотах.

\begin{center}
\begin{tabular}{l|cccc|c}
привод 1 & $0{,}110$~Гц & $0{,}391$~Гц & $0{,}439$~Гц & $0{,}635$~Гц & разброс\\
\hline
$F_c$, Н$\cdot$м & $61{,}0$ & $79{,}1$ & $97{,}9$ & $80{,}0$ & $16\%$\\
$C_v$, Н$\cdot$м$\cdot$с & $531$ & $238$ & $183$ & $190$ & $50\%$\\
$\tau$, мс & $1054$ & $265$ & $206$ & $156$ & $88\%$\\
\end{tabular}
\end{center}

Транспортная задержка исключается: чтобы объяснить невязку, $\tau$ должна меняться в семь раз монотонно с частотой. Фазовый сдвиг первой гармоники составляет $29\ldots46^{\circ}$ и практически не зависит от частоты в диапазоне $0{,}11\ldots0{,}64$~Гц, что несовместимо с задержкой в канале измерения, но характерно для сухого трения: его момент примерно постоянен и входит в квадратуре к основному, гравитационному, который от частоты также почти не зависит.

Оценка момента потерь составляет $48\ldots67$~Н$\cdot$м на выходном валу. Независимая проверка по паспорту мотор-редуктора: при передаточном числе $i = 81{,}5$, номинальном моменте двигателя $M_1 = 10{,}2$~Н$\cdot$м и моменте на выходном валу $M_2 = 773$~Н$\cdot$м КПД равен $M_2/(M_1 i) = 0{,}930$, что соответствует моменту потерь $M_1 i - M_2 = 58{,}3$~Н$\cdot$м. Собственные механические потери двигателя добавляют ещё $8\ldots17$~Н$\cdot$м, приведённых к выходному валу.

\subsection{Проверка в статике}

Модель \ref{eq:FrictionModel} в статическом режиме даёт три дополнительных неизвестных $F_{c,j}$, $j=1,2,3$, входящих линейно; знаки $\operatorname{sign}\dot{q}$ определяются по направлению подхода к стационарному положению и считываются из записи углов.

Оценивание даёт величины не более $20$~Н$\cdot$м, причём знаки меняются при переходе от $\gamma = 0{,}1$ к большому $\gamma$, а доверительные интервалы в большинстве случаев накрывают нуль. Статическое трение по имеющимся данным не идентифицируется, и это физически объяснимо: при остановке сила трения не остаётся на максимуме, а принимает любое значение внутри зоны застоя, необходимое для равновесия.

Практическое следствие: монотонный подход ко всем стационарным положениям не внёс заметного смещения в оценки. Добавление трения в статическую модель сдвигает $x_c$, $y_c$, $f_0$ меньше, чем на ширину их доверительных интервалов, и при этом расширяет интервалы для $k_s$ и $z_c$. В статическую модель трение включать не следует.

\section{Приведённая инерция ротора}

Допущение о пренебрежимости массоинерционных характеристик вспомогательных звеньев не распространяется на ротор двигателя: при передаточном числе $i = 81{,}5$ приведение к валу кривошипа идёт с множителем $i^2 = 6642$. Для двигателя мощностью $1{,}5$~кВт с тормозом $J_m = 32\cdot10^{-4}$~кг$\cdot$м$^2$, что даёт
\begin{equation}
J_{\text{прив}} = J_m i^2 \approx 21{,}3~\text{кг}\cdot\text{м}^2 .
\end{equation}

Вклад в момент на валу растёт как $f^2$:

\begin{center}
\begin{tabular}{l|cccc}
частота, Гц & $0{,}110$ & $0{,}391$ & $0{,}439$ & $0{,}635$\\
\hline
$J_{\text{прив}}\ddot{q}$, Н$\cdot$м & $2{,}2$ & $21{,}4$ & $44{,}2$ & $57{,}4$\\
доля от расчётной амплитуды & $2{,}5\%$ & $15\%$ & $19\%$ & $29\%$\\
\end{tabular}
\end{center}

На гармонических участках движение практически одностепенное, поэтому $\ddot{q}_i$ пропорционально $\ddot{\alpha}$, и слагаемое $J_{\text{прив}}\ddot{q}$ почти коллинеарно слагаемому $c_xI_x\ddot{\alpha}\cos{\beta}$. Разделить их по имеющимся данным нельзя: подгонка коэффициентов $c_x$, $c_y$ уже впитала приведённую инерцию ротора. Следовательно, настроенные $c_x$ и $c_y$ оценивают не моменты инерции кабины, а их сумму с приведённой инерцией привода.

Величина $J_{\text{прив}}$ известна из паспорта и не подлежит оцениванию: её место -- в известной части модели. Вычитание $J_{\text{прив}}\ddot{q}$ из измеренного момента до настройки $c_x$, $c_y$ устраняет систематическую погрешность до $29\%$ и не требует новых экспериментов.

\section{Результаты оценивания}

Доверительные интервалы получены ресэмплингом стационарных участков (бутстрэп, 200 выборок, уровень $90\%$). В интервалы включён разброс по вариантам метода: $\gamma = 0{,}1$ и $\gamma \to \infty$, с учётом трения и без него.

\begin{center}
\begin{tabular}{l|c|c}
параметр & оценка & $90\%$ интервал\\
\hline
$f_0$, Н -- преднатяг пружины & $14\,700$ & $14\,500\ldots14\,800$\\
$z_c$, м -- высота центра тяжести & $0{,}87$ & $0{,}84\ldots0{,}91$\\
$y_c$, см -- поперечное смещение & $-6{,}8$ & $-7{,}15\ldots-6{,}21$\\
$x_c$, см -- продольное смещение & $-0{,}67$ & $-0{,}79\ldots-0{,}52$\\
$k_s$, Н/м -- жёсткость пружины & $\approx13\,000$ & $11\,000\ldots18\,400$\\
\end{tabular}
\end{center}

Пружина несёт $14{,}7$~кН из $20{,}1$~кН веса кабины, то есть $73\%$.

Невязка по моментам при $\gamma = 0{,}1$: среднее $17$~Н$\cdot$м, СКО $23$~Н$\cdot$м, максимум $66$~Н$\cdot$м при рабочих моментах $200\ldots500$~Н$\cdot$м. Эта величина и служит показателем качества модели в статике.

Слабо определённым остаётся лишь коэффициент жёсткости ($\pm30\%$). Причина структурная: $k_s$ входит только через произведение $k_s\bar{h}/L$, а участки с большим вертикальным ходом (номера $1-10$, $h$ от $-0{,}145$ до $+0{,}033$~м) непригодны из-за мёртвых точек. Для уточнения жёсткости требуется отдельная серия с вертикальным перемещением при углах кривошипов вдали от нуля.

\section{Замечания к расчётным реализациям}

При сверке MATLAB- и Python-версий выявлены три расхождения.

\textbf{Коэффициент жёсткости.} В \texttt{calc\_Jmin\_zc.m} коэффициентом при $k_s$ служит $\bar{h}$, тогда как уравнение сил требует $\bar{h}/L$. Поэтому соответствующая компонента вектора $\tau$ равна $k_s/L$, а не $k_s$: опубликованное значение $24\,606$ отвечает истинной жёсткости $24\,606\cdot L \approx 13\,800$~Н/м. Дополнительно, при расчёте правой части $B$ слагаемое с $k_s$ делится на $L$ повторно.

\textbf{Коэффициент $C_{2,h}$.} В \texttt{KinematicDirProb.m} и в настоящем выводе коэффициент при $\sin^2{q_3}$ равен $a$; в \texttt{ident\_parameters.m} в этом месте стоит $2a$. Проверка алгеброй подтверждает $a$. Расхождение влияет на поправки порядка $\kappa^3$ к $S^0_x$ и $S^0_y$.

\textbf{Перекрёстные инерционные члены.} Выражения $\left(I_y-I_x-I_z\right)\dot{\alpha}\dot{\beta}\sin{\beta}$ и $\left(I_z-I_x\right)\dot{\alpha}^2\sin{\beta}\cos{\beta}$ в расчётных реализациях опущены. На имеющихся данных это оправдано: $\left|\beta\right| \le 3^{\circ}$, вклад слагаемых порядка $1$~Н$\cdot$м против сотен, -- но в изложении отбрасывание следует оговаривать.
